% Russian protocol and provenance supplement.
\documentclass[11pt,a4paper]{article}
\usepackage[T2A]{fontenc}
\usepackage[utf8]{inputenc}
\usepackage[english,russian]{babel}
\usepackage[margin=25mm]{geometry}
\usepackage{amsmath,amssymb,booktabs,longtable,array}
\usepackage[numbers,sort&compress]{natbib}
\usepackage[hyphens]{url}
\usepackage[hidelinks]{hyperref}
\ifx\pdfmapfile\undefined\else
\pdfmapfile{+cm-super-t1.map}\pdfmapfile{+cm-super-t2a.map}
\pdfmapfile{+cm-super-ts1.map}\fi
\setlength{\emergencystretch}{2em}
\renewcommand{\thesection}{S\arabic{section}}
\renewcommand{\theequation}{S\arabic{equation}}
\renewcommand{\thetable}{S\arabic{table}}

\title{Дополнительные материалы\\
Передача напряжений и движение дислокаций в алюминии
вблизи включения Al$_{13}$Fe$_4$ с удерживаемой деформацией}
\author{И.\,Е.~Михайловский \and Д.\,Е.~Пшонкин}
\date{Московский политехнический университет, Москва, Россия\\
\texttt{ilyamihailovsy@gmail.com}}

\begin{document}
\maketitle

\section{Структура и подготовка}

Ячейка границы содержит 91\,428 атомов.
Идентификаторы 1--72\,532 принадлежат матрице,
остальные 18\,896 атомов относятся к включению.
Входное соответствие типов LAMMPS задано в порядке
Al, Fe, Mg и Si, хотя в конфигурациях границы Mg
и Si отсутствуют. Полная ячейка содержит
87\,076 атомов Al и 4352 атома Fe.
Оси ГЦК-алюминия совпадают с направлениями
$[1\bar{1}0]$, $[11\bar{2}]$ и $[111]$,
параметр решётки равен 4{,}05~\AA{}.
Al$_{13}$Fe$_4$ имеет моноклинную структуру
\cite{Liu2024CIF}. Один потенциал Елинека
используется внутри фаз и на границе \cite{Jelinek2012}.

Исходные размеры ячейки равны
$L_x=154{,}644253$~\AA{},
$L_y=64{,}482817$~\AA{}; границы по $z$
составляют $-10{,}0$ и 165{,}943041~\AA{}.
При построении матрицы используются 56 слоёв (111).
Верхняя поверхность материала отличается от границы
расчётной области и смещается при подготовке,
поэтому единственная точная высота свободной
поверхности для всех состояний не задаётся.
Опорный слой доходит до $z=20$~\AA{}.
Выступ центрирован при $x=77{,}3221$~\AA{}
и имеет полуоси 35 и 20~\AA{} в плоскости $xz$.

Для согласования периодичности включение
масштабируется в плоскости с коэффициентами
0{,}9978336 и 0{,}9973955.
Атомы матрицы, расположенные ближе 2{,}2~\AA{}
к атомам включения, удаляются при построении.
Минимальное расстояние после этой операции
составляет примерно 2{,}219~\AA{} внутри ячейки
и 2{,}292~\AA{} через периодический шов.
Геометрическая проверка исключает исходные
перекрытия на данном масштабе расстояний,
но не проверяет энергию и структуру
релаксированной границы. Область $z<6$~\AA{}
закреплена.

В LAMMPS используются единицы \texttt{metal},
атомный тип частиц и периодические границы
вдоль $x$ и $y$ при фиксированных границах
вдоль $z$ \cite{Thompson2022LAMMPS}.
Запас списка соседей равен 2~\AA{}.
При минимизации список обновляется на каждом
шаге. Подготовка контроля начинается
с минимизации методом сопряжённых градиентов
при заданной евклидовой норме сил 1~эВ/\AA{}.
Затем выполняются 6~пс при 300~К
с термостатом Нозе--Гувера и временем
демпфирования 0{,}1~пс. Ланжевеновское
охлаждение снижает целевую температуру
до 0{,}5~К за 12~пс при времени
демпфирования 0{,}2~пс.
Начальные значения генераторов случайных
чисел для скоростей и охлаждения равны
4242 и 4243. Перед заключительной минимизацией
скорости обнуляются; заданная норма сил
составляет 0{,}05~эВ/\AA{}.
В ходе тепловой подготовки у верхней границы
ячейки действует отражающая стенка.
Полученное состояние служит исходным
для парного статического расчёта.

Обе конфигурации статической пары
восстанавливаются из одного контрольного
снимка. Смещаемая часть выступа содержит
4816 атомов включения при $z>20$~\AA{}.
Сохранённые координаты её центра
примерно равны
$(77{,}258,32{,}201,28{,}678)$~\AA{}.
Тензор возмущения имеет вид
\begin{equation}
\mathbf E_0=
\begin{pmatrix}
0.000485&0&0.001455\\
0&-0.000970&0\\
0.001455&0&0.000485
\end{pmatrix}.
\label{eq:tensor}
\end{equation}
Опорная координата меняется на
$\delta\mathbf r=\mathbf E_0(\mathbf r-\mathbf r_c)$.
Наибольшее заданное смещение равно
0{,}05412~\AA{}.
Определитель $\mathbf I+\mathbf E_0$ равен
$(1+\varepsilon_0)(1-\varepsilon_0/2)^2$
при $\varepsilon_0=0{,}00194$.
Его отклонение от единицы составляет
около $2{,}82\cdot10^{-6}$, что соответствует
сохранению объёма только в первом порядке
по деформации.

Все атомы включения удерживаются потенциалом
\begin{equation}
U_{\mathrm s}=\frac{K}{2}\sum_{i\in\mathrm{inclusion}}
 |\mathbf r_i-\mathbf r_i^{\mathrm{ref}}|^2,
\qquad K=20~\text{эВ}/\text{\AA}^2 .
\label{eq:spring}
\end{equation}
Опорные координаты нижнего слоя не смещены
в обеих копиях; координаты выступа
различаются согласно уравнению~\eqref{eq:tensor}.
Энергия пружин включается в минимизацию.
Для сопряжённых градиентов заданы
максимальное смещение 0{,}05~\AA{},
нулевой энергетический допуск,
допуск по норме сил 0{,}02~эВ/\AA{},
не более 6000 итераций и 60\,000 вычислений
сил. Контроль с удерживаемым включением
останавливается по силовому критерию
после 235 итераций при конечной норме
0{,}019779915~эВ/\AA{}.
Возмущённое удерживаемое состояние также
достигает силового критерия после
310 итераций при
0{,}019879873~эВ/\AA{}.
Эти значения подтверждают выполнение
заданного условия остановки.
Сходимость профиля напряжений по допуску,
геометрии и начальному состоянию
ими не устанавливается.

\section{Напряжения и расстояния}

Команда \texttt{compute stress/atom}
с параметрами \texttt{NULL virial} сохраняет шесть компонент
в порядке $xx,yy,zz,xy,xz,yz$.
Для выбранного по контролю набора из
$N$ атомов Al используется пересчёт
\begin{equation}
\sigma_{\alpha\beta}
 =-\frac{1}{10N\Omega}\sum_{i=1}^{N}S_{i,\alpha\beta},
\qquad \Omega=16{,}6072~\text{\AA}^3,
\label{eq:stress}
\end{equation}
где $S_i$ записаны в бар~\AA$^3$,
а результат выражается в МПа.
При таком пересчёте нормальные напряжения положительны при сжатии,
а тензор противоположен по знаку тензору Коши с положительным
растяжением. Это соглашение сохранено для всех величин МД со знаком.
При сравнении с континуальной моделью используются модули проекций,
не зависящие от общей смены знака. Кинетический вклад
не включён. Полученная конфигурационная
характеристика не является измерением
температуры или теплового потока.

Атомные идентификаторы и типы элементов
сопоставляются между снимками.
Координаты контроля определяют
принадлежность слою и осевому окну
шириной 20~\AA{}.
Центр окна определяется серединой
между крайними координатами атомов
контроля по $x$.
Слои толщиной 4~\AA{} по $z$
содержат атомы типа 1.
Учитываются только слои, нижняя граница которых выше всех атомов
включения в обоих снимках. Включению принадлежат атомы Al и Fe
с идентификаторами больше 72\,532. Их максимальная высота равна
39{,}707546~\AA{} в контроле и 39{,}718591~\AA{} в возмущённом
состоянии. Центры 23 учитываемых слоёв находятся между
$r=22$ и 110~\AA{}, где
$r=z-20$~\AA{}.
Последний слой ограничен
$r=108$ и 112~\AA{}.
Такой отбор исключает слои с алюминием
включения и самые верхние поверхностные
слои.

Тензор усредняется отдельно в каждой
конфигурации и затем вычитается.
Для
$\mathbf D=\Delta\boldsymbol\sigma-
\operatorname{tr}(\Delta\boldsymbol\sigma)\mathbf I/3$
напряжение фон Мизеса равно
$\sqrt{3\mathbf D:\mathbf D/2}$.
Максимальная по модулю проекция
на систему скольжения определяется
как
$\max_\alpha|\mathbf s_\alpha\cdot
\Delta\boldsymbol\sigma\,\mathbf n_\alpha|$.
Компонента $xz$ и её знак сохраняются
отдельно. Поэтому малое среднее
касательное напряжение по ширине
может сочетаться с большим
напряжением противоположного знака
в осевом окне.

Статистика удалённых слоёв рассчитана
по тринадцати слоям при $r>60$~\AA{}.
Средний максимум проекции равен
2{,}477~МПа, пространственное
стандартное отклонение составляет
0{,}543~МПа.
Слои не являются независимыми
реализациями однородного состояния.
Эти характеристики описывают
пространственную неоднородность
и не задают стандартную ошибку,
численный шум или физический
порог обнаружения.
Остаточное напряжение на наибольшем
исследованном расстоянии
указывается непосредственно.

Высота ослабления определяется
только относительно заданного
уровня напряжений.
Если $h_m(s)$ обозначает высоту
над номинальной вершиной,
на которой убывающий модуль
осевого сдвига достигает $s$,
линейная интерполяция соседних
слоёв даёт примерно
4{,}69, 6{,}69 и 8{,}55~нм
для $s=10$, 5 и 2~МПа.
Это операционные характеристики
профиля с шагом 0{,}4~нм,
а не независимо установленные
длины материала.
Высота полного исчезновения
напряжения не определяется.

\section{Сохранённая деформация и свободные диагностические расчёты}

Методом наименьших квадратов подбирается
аффинное отображение между положениями
соответствующих атомов Fe
в двух минимизированных конфигурациях.
После исключения переноса получается
градиент деформации $\mathbf F$.
В периодических направлениях
используются минимальные
межобразные смещения.
Сохранённая деформация задаётся
\begin{equation}
\mathbf E_{\mathrm{res}}=
 \tfrac12(\mathbf F^{\mathsf T}\mathbf F-\mathbf I),
\qquad
\eta=\frac{\mathbf E_{\mathrm{res}}:\mathbf E_0}
                 {\mathbf E_0:\mathbf E_0}.
\label{eq:eta}
\end{equation}
Для 1008 атомов Fe с контрольной
координатой $z>22$~\AA{}
получается $\eta=0{,}9728$.
Основные компоненты равны
$E_{xx}=0{,}000482$,
$E_{yy}=-0{,}000887$,
$E_{zz}=0{,}000460$ и
$E_{xz}=0{,}001435$.
Среднеквадратичный остаток подгонки
составляет 0{,}0035~\AA{}.

Формальная стандартная ошибка
0{,}0038 получена распространением
ковариации коэффициентов
метода наименьших квадратов
при объединённой дисперсии
остатков и $3N-12$ степенях
свободы.
Атомные остатки пространственно
коррелированы; эта характеристика
не включает ошибку модели
или разброс независимых
траекторий.
Скалярная проекция также
не означает сохранения
одинаковой доли каждой
компоненты тензора.

Минимизации без удержания
сохраняются только как
диагностические данные.
Для возмущённой пары без
пружин в сохранённом результате
проекция выступа равна
$0{,}4364\pm0{,}1019$.
Обе минимизации остановлены
после 6000 итераций
при конечных нормах сил
0{,}8660473 и
0{,}48161949~эВ/\AA{}
для контроля и возмущённого
состояния соответственно,
выше заданных
0{,}02~эВ/\AA{}.
Отдельная запись прямой
минимизации даёт
$0{,}2007\pm0{,}1064$
для иначе подготовленной
пары.
Ни один из результатов
не устанавливает равновесную
сохранённую долю, время
релаксации или поле после
снятия магнитного воздействия.
Эти профили не используются
в основном количественном
сопоставлении.

\section{Нагружение и отслеживание линий}

Нагружаемая ячейка содержит
91\,428 атомов и сохраняет
ориентацию и форму включения
статического расчёта.
Центры исходных нижней
и верхней краевых линий
приблизительно равны
$(139{,}322,49{,}228)$ и
$(115{,}922,72{,}611)$~\AA{}
в плоскости $xz$.
Смещения задаются
периодической конструкцией
Вольтерры.
Координаты введения
отличаются от координат
после минимизации
и начальной динамики.

Входной файл нагружения
сначала минимизирует
построенную ячейку
с заданной нормой сил
1~эВ/\AA{} и пределом
2000 итераций.
Затем задаются скорости
при 300~К с начальным
значением генератора
88004; шаг интегрирования
равен 1~фс.
Нижняя область закреплена,
подвижные атомы связаны
с термостатом Нозе--Гувера
с временем демпфирования
0{,}1~пс, у верхней
границы ячейки действует
отражающая стенка.
Одинаковые силы вдоль
$x$ на выбранных верхних
атомах передают номинальную
сдвиговую нагрузку.
В зависимости от варианта
включение свободно
либо связано с пружинами
жёсткостью 20~эВ/\AA$^2$.

Пусть $s$ обозначает
число шагов после
начальной выдержки
длительностью 5000 шагов,
$T_s=16\,000$ шагов,
а $N$ является полной
длиной расчёта.
Номинальное напряжение
в МПа задаётся выражением
\begin{equation}
\tau(s)=
\begin{cases}
0,&s<0,\\
q\left[\dfrac{s}{2}-\dfrac{T_s}{2\pi}
 \sin\left(\dfrac{\pi s}{T_s}\right)\right],
 &0\le s<T_s,\\
q(s-T_s/2),&s\ge T_s,
\end{cases}
\qquad
q=\frac{\tau_{\max}}{N-5000-T_s/2}.
\label{eq:ramp}
\end{equation}
Для свободного контроля
$N=101\,000$,
$\tau_{\max}=400$~МПа.
Для удерживаемых контроля
и возмущённой ячейки
$N=45\,000$,
$\tau_{\max}=145{,}45$~МПа.
Скорости после плавного
начала совпадают
с точностью указанных
округлённых параметров.
Регулярные снимки
сохраняются через
2000 шагов.
В таблицах анализа использовано заданное конечное напряжение
каждого расчёта. Разность между поздними кадрами равна
9{,}0909~МПа в свободном контроле и 9{,}090625~МПа
в удерживаемой паре.

DXA определяет
векторы Бюргерса
и сегменты линий
\cite{Stukowski2012DXA,Stukowski2010OVITO}.
Для положений по $x$
снимается периодическая
неоднозначность при
$L_x=154{,}644253$~\AA{}.
Для каждого семейства
вектора Бюргерса
линейный базовый тренд
подбирается по шагам
не ранее 4000 при
номинальном напряжении
ниже 40~МПа.
Критерий ухода задаётся
модулем остатка
больше
$\max(6s_x,8~\text{\AA})$,
где $s_x$ является
стандартным отклонением
остатков базового тренда.
Проверяются следующие
две позиции наблюдения,
если они доступны;
отсутствие линии
допускается сохранённым
детектором.
Численный критерий
поэтому требует
отдельной проверки
непрерывности у границ
и после исчезновения
сегмента.

Для верхней линии
стандартные отклонения
остатков базового
тренда равны
3{,}8919 и
3{,}5046~\AA{}
в удерживаемом контроле
и возмущённой ячейке.
Соответствующие пороги
составляют 23{,}351
и 21{,}027~\AA{}.
Первые принятые
кадры относятся
к 38 и 40~пс.
Таким образом,
одна процедура
использует разные
подобранные тренды
и пороги смещения.
Разница в один
кадр не является
независимым измерением
упрочнения на
9{,}09~МПа.

Нижнее семейство
исчезает к 10~пс
в обоих удерживаемых
расчётах.
Его последующие
координаты могут
относиться к сегментам
шва и не задают
второй достоверный
уход исходной линии.
В последнем
проанализированном
кадре контроля
при 44~пс суммарная
извлечённая длина
составляет
39{,}344~\AA{}
в трёх сегментах.
Длины, суммированные
по семействам,
равны 26{,}191
и 13{,}153~\AA{};
они не обязательно
совпадают с длинами
отдельных связных
линий.
В соответствующем
возмущённом кадре
сегментов не выделено.
Отсутствие обнаруженной
линии отделяется
от установления
атомистического пути
реакции.

В контроле со
свободным включением
уход нижней линии
зарегистрирован
при 36~пс
и 104{,}545~МПа;
предыдущий кадр
соответствует
95{,}455~МПа.
Критерий верхней
линии выполняется
при 40~пс
и 122{,}727~МПа.
Последний регулярный
кадр анализа
относится к 100~пс
и 395{,}455~МПа;
номинальные
400~МПа таблицей
не охвачены.
Короткие сигналы
на шве встречаются
при 232--241~МПа
и повторно выше
377~МПа.
Они не подтверждают
устойчивого зарождения
на границе.
Отсутствие подтверждённого
события не позволяет
приписать скалярному
порогу значение
выше 395{,}455~МПа.

\section{Отдельная ячейка сплава}

Ячейка сплава содержит
92\,800 атомов.
Размеры в плоскости
равны
$114{,}5513\times49{,}6022$~\AA{},
границы по $z$
составляют $-8{,}0$
и 283{,}2392~\AA{}.
В ней находятся
91\,315 атомов Al,
928 Mg и 557 Si.
Номинальные
концентрации
по полной ячейке
равны 1{,}0~ат.\% Mg
и примерно
0{,}6~ат.\% Si.
Замещения ограничены
подвижной областью;
граничные слои
остаются алюминиевыми.
Нижний закреплённый
слой имеет толщину
8~\AA{}, верхний
нагружаемый
12~\AA{}.
Расстояние между
краевыми линиями
по нормали равно
205{,}768~\AA{}.

Начальная минимизация остановилась по допуску энергии при
норме сил 10{,}118429~эВ/\AA{} и максимальной компоненте силы
1{,}9509711~эВ/\AA{}. Заданный допуск сил не достигнут.
Последующая траектория при 300~К завершена, однако такая подготовка
не подтверждает механическое равновесие исходного расположения
дислокаций. Траектория содержит
10~пс без
приложенного сдвига
и последовательные
номинальные уровни
45, 55, 65
и 75~МПа
по 30~пс.
Каждое увеличение
сглажено за 4~пс.
Суммарное смещение нижней линии от 10 до 130~пс равно
1{,}98~\AA{} при стандартном отклонении положения 0{,}63~\AA{}.
Начальная выдержка без нагрузки сюда не входит. Смещение
по всей траектории от 0 до 130~пс составляет 6{,}99~\AA{}.
Смещения на
отдельных уровнях
составляют
0{,}58, 1{,}56,
$-1{,}53$ и
0{,}17~\AA{}.
Смещение за установившийся участок каждой ступени меньше
модуля вектора Бюргерса в этой реализации расположения примесей.

Наблюдения зависят
от выбранного
расположения примесей,
внутренних напряжений
и короткого
времени наблюдения.
Приложенное
напряжение
не измеряет
локальную
разрешённую
касательную нагрузку
на дислокации.
Траектории
подтверждают
ограниченное движение
при заданной программе,
но не устанавливают
собственную прочность
препятствий
75~МПа.
Активационный объём
по подобранным
наклонам
не определяется.

Открытый набор содержит полную сохранённую атомную траекторию
с исходными столбцами координат и напряжений в сжатых частях
\path{data/publication/alloy/trajectory.*.lammpstrj.gz}.
Также представлены производная таблица положений из 131 кадра,
сводные результаты, начальная и конечная структуры и журнал.
Доступный начальный файл сохранён после даты расчёта. Идентификаторы,
типы атомов и координаты закреплённой границы согласуются с траекторией,
но исходной копии или контрольной суммы на момент запуска нет.
Совпадение координат подвижных атомов поэтому не подтверждено.
Этот файл служит возможной стартовой структурой, без гарантии точного
повторения исходного запуска. Сохранённая траектория позволяет повторить DXA и проверить
пространственный отбор, использованный при анализе сплава.
Активационный объём из этих наблюдений не определяется.

При первоначальном усреднении напряжений в выбранной области учитывались
Al и Si, а Mg (тип 2) исключался. Повторный анализ всех 131 сохранённого
кадра воспроизводит прежние смещения и напряжения. Отбор только Al или
всех атомов сплава меняет среднее напряжение на ступени менее чем
на 1{,}71~МПа. Диагностические оценки $V^*/b^3$ равны
$-12{,}32\pm21{,}31$ для исходного отбора, $-13{,}54\pm19{,}19$
для Al и $-12{,}32\pm19{,}83$ для всех атомов при положительном сжатии.
Переход к положительному растяжению меняет знаки этих оценок.
Подгонка колеблющихся
положений не определяет физический активационный объём. Сравнение
сохранено в \path{stageG6_stress_selection_relA.json}.

\section{Континуальная и активационная оценки}

Сферическая модель
имеет тензор
уравнения~\eqref{eq:tensor},
радиус 35~\AA{},
модуль сдвига
26{,}5~ГПа
и коэффициент
Пуассона 0{,}347.
Изотропные
упругие постоянные
включения
и бесконечной
матрицы одинаковы.
Через тензор
Эшелби $\mathbf S$
и жёсткость
$\mathbf C$
внутреннее
напряжение
выражается как
\begin{equation}
\boldsymbol\sigma_{\mathrm{in}}
 =\mathbf C:(\mathbf S:\mathbf E_0-\mathbf E_0).
\label{eq:eshelby}
\end{equation}
Максимальная
внутренняя проекция
на систему скольжения
равна 41{,}45~МПа
\cite{Eshelby1957}.
Это внутреннее
значение в задаче
с собственной
деформацией,
которое не задаёт
поверхностную амплитуду
каждой внешней
компоненты.
У выступа
основного расчёта
иные ограничения,
опорный слой
и близкая свободная
поверхность.
Сопоставление
не устанавливает
порядка величин
полей для всех
условий и
не проверяет
атомистическое
граничное условие.

Скалярный расчёт
ползучести принимает
внешнюю амплитуду
компактной частицы
$\tau_m(a/R)^3$
и объёмную долю
$f$.
В первом порядке
по $f$
пространственный вес
модулей напряжения
между $s$
и $s+\mathrm ds$
равен
$f\tau_m s^{-2}\mathrm ds$
при
$0<s<\tau_m$.
Этот вес используется
только в интеграле
добавочной скорости.
Нормированной
плотностью вероятности,
продолженной
до нулевого
напряжения,
он не является.
Равные веса
знаков напряжения
и независимый
от напряжения
активационный объём
дают
\begin{equation}
H=f x\int_0^x\frac{\cosh u-1}{u^2}\,\mathrm du
 =f\{x\,\mathrm{Shi}(x)-[\cosh(x)-1]\},
\qquad
x=\frac{V^*\tau_m}{k_{\mathrm B}T}.
\label{eq:rate}
\end{equation}
Здесь
$\mathrm{Shi}(x)=\int_0^x\sinh(u)/u\,\mathrm du$.
Подынтегральная
функция имеет
конечный предел
$1/2$ в нуле.
Перекрытие частиц,
распределение
ориентаций зёрен
и зависимость
формы барьера
от напряжения
не учитываются.

В оценке
чувствительности
использованы
$k_{\mathrm B}T=4{,}141947\cdot10^{-21}$~Дж
при 300~К,
$b=2{,}86378246$~\AA{}
и $f=0{,}00246$.
Объёмная доля
представляет
модельный сценарий.
Для массовой
доли $w=0{,}0035$
и плотностей
$\rho_i=3{,}85$,
$\rho_m=2{,}70$~г/см$^3$
получается
\begin{equation}
f=\frac{w/\rho_i}{w/\rho_i+(1-w)/\rho_m}
 \simeq0{,}00246 .
\end{equation}
Принятые плотности
и состав
не определяют
объёмную долю
каждого
экспериментального
образца.
Диапазон
$19b^3$--$142b^3$
используется
для анализа
чувствительности.
Опубликованное
значение $70b^3$
для Al--Mg--Si
после равноканального
углового прессования
характеризует
другое состояние
материала
\cite{Soula2022JALCOM}.

При $H=0{,}25$
уравнение~\eqref{eq:rate}
даёт
$x=6{,}754241794$.
Необходимые
амплитуды равны
62{,}691,
39{,}705,
23{,}823,
17{,}016,
11{,}911
и 8{,}388~МПа
при $V^*/b^3=19$,
30, 50, 70,
100 и 142.
При фиксированных
$\tau_m=15{,}451$
и 41{,}45~МПа
получаются
$V^*/b^3=77{,}091$
и $28{,}737$.
Первая амплитуда
является максимумом
в осевом окне
выступа,
вторая внутренней
проекцией сферы.
Их подстановка
в скалярную модель
внешнего поля
явно условна.
Она не калибрует
магнитострикцию
и не устанавливает
экспериментальную
скорость ползучести.

\section{Записи и воспроизводимость}

Публикационные
данные и код
размещены по адресу
\url{https://github.com/iluuua/magnetostriction-plasticity-bounds}.
Манифест в
\path{data/publication}
описывает исходные
конфигурации,
сохранённые столбцы
и контрольные суммы.
В нём разделены
полные снимки
атомных вириалов
статической пары,
траектория сплава со всеми исходными столбцами
и координатные проекции трёх траекторий нагружения границы.
Проекции сохраняют
идентификаторы,
типы, координаты,
времена кадров
и границы ячейки
для каждого
записанного
регулярного кадра.
Все исходные
поатомные столбцы
напряжений
и энергий
в них не включены.

Напряжения
основаны на
\path{stageG10_field_profile.json}
и соответствующем
CSV, рассчитанных
по удерживаемым
контрольному
и возмущённому
снимкам с вириалами.
Отбор слоёв,
проекции со знаком
и максимумы
можно пересчитать
по первичным
статическим данным.
Компоненты
сохранённой
деформации
и формальная
ошибка подгонки
приведены в
\path{stageG12_eigenstrain_retention.json}.
Свободные
диагностические
результаты
содержатся в
\path{stageG12_eigenstrain_retention_free.json}
и
\path{stageG12_eigenstrain_retention_free_cg.json}.
Их использование
ограничено отсутствием
сходимости,
как указано
в разделе~S3.

Движение вблизи
границы описано
в
\path{stageG2_depinning_summary_G15held.json}
и
\path{stageG2_depinning_summary_G15ctl_free.json}.
Соответствующие
таблицы
\path{stageG2_depinning_G15_ctl_G15held.csv},
\path{stageG2_depinning_G15_fld_G15held.csv}
и
\path{stageG2_depinning_G15_ctl_G15ctl_free.csv}
содержат времена
кадров,
номинальные
напряжения,
положения
и длины
по семействам.
Опубликованные
координатные
траектории позволяют
независимо
восстановить
линейные дефекты.
Сводки детектора
требуют проверки
непрерывности
из раздела~S4;
поздняя запись
семейства
не обязательно
относится
к исходной линии.

Для ячейки сплава
использованы
\path{stageG6_vstar_relA_frames.csv} и
\path{stageG6_vstar_relA_summary.json}.
В работу включены
только история
положений
и завершённый
диапазон нагрузок.
Подобранный
активационный объём
из диагностической
сводки физической
оценкой не считается.
Внутреннее поле
сферы соответствует
\path{stageG8_eshelby3d.json}.
Интеграл
и обратные
оценки напряжения
определяются
уравнением~\eqref{eq:rate}
с указанными
выше параметрами;
соответствующий
архивный расчёт
содержится в
\path{stageG5_two_scale_bridge.json}.
Текстовые трактовки
в записях
не заменяют
эти определения
и сведения
о сходимости.

Код построения
рисунков работает
с указанными
конфигурациями
и записями.
Изображения моделей
созданы в OVITO
и представляют
рассчитанные
атомные координаты.
Они не являются
экспериментальными
микрофотографиями.
Четыре основных
рисунка показывают
геометрию,
статическое поле,
траектории
вблизи границы
и условную
активационную
зависимость.
Подписи выполнены
на языке рукописи.
Численные данные
на рисунках
обеих версий
одинаковы.

\bibliographystyle{unsrtnat}
\bibliography{references}
\end{document}
